\section{Introduction}\label{sec:introduction}

\paragraph{The question.}
Let $b(d)=\max(1,\lceil\log_2d\rceil)$ and $L=d^2b(d)$. Question~(T) asks whether one
deterministic multitape Turing machine, with a fixed finite number of tapes and a
fixed finite alphabet, transposes every $d\times d$ matrix with $b(d)$-bit entries,
given and returned in row-major order, in $O(L)$ steps.

The source-cover paper of this series~\cite{SourceCovers2026} fixes one complete
computation geometry and a nonempty set $\mathcal C$ of inputs that it transposes
correctly. It bounds the correctness deficit $\Delta=L-\log_2|\mathcal C|$ below by a
geometric credit minus the information revealed by a certificate
\cite[Theorem~3.5]{SourceCovers2026}. At each legal phase $\rho$ and column $j$ the
credit is at most the circular point-gap area $\gaparea(S_\rho(s_j))$ of the complete
raw-source set of an exact source bank, frozen at the first output time $s_j$ of the
column \cite[Theorems~2.7 and~2.9]{SourceCovers2026}. That paper bounds the frozen area
below by fractional covers, anchored approximation and shared support. It obtains a
uniform exclusion for one exceptional family, and further exclusions under
economical-cover or shared-support conditions. For two exceptional families it
identifies the inverse-square hop profiles as the nearest unresolved case
\cite[Section~9.1]{SourceCovers2026}.

\paragraph{Results.}
\begin{enumerate}
\item \emph{Theorem~G (\cref{thm:necessary-gap-area}).} We keep the frozen point area
itself as the certificate. Along every joint profile of every near-full sequence, and
for all integers $R\ge0$ and $M\ge\max(2,R+1)$,
\[
 \limarea_{R,M}\le\frac{\beta\tau_R}{M-R},
\]
where $\limarea_{R,M}$ is the limiting mean frozen point area over legal phases and
columns, with the first $R$ core families protected (\cref{def:gap-area-means}). The
limits are those of the reciprocal-cover bound \cite[Theorem~4.6]{SourceCovers2026},
taken in the same order. Since the frozen point area dominates the fractional-cover
credit \cite[Lemma~3.3]{SourceCovers2026}, Theorem~G implies that bound
(\cref{rem:reciprocal-comparison}).
\item \emph{Theorem~A (\cref{thm:funnels}).} A funnel is a property of the full family
system: time sets $D_0,\dots,D_m$ such that each $D_i$ with $i<m$ is fed by at most
one family, which moves into $D_i$ only times of $D_{i+1}$, while no family feeds
$D_m$ (\cref{def:funnel}). Every orbit then meets $D_0$ in at most $m+1$ times
(\cref{lem:funnel-area}). With Theorem~G, no near-full sequence has, from some index on
or along a subsequence, funnels of a fixed depth whose first set contains every time
that reads a raw interval of fixed positive relative length. The number of exceptional
families and the sizes of the core families are arbitrary.
\item \emph{Theorem~B (\cref{thm:terminal-landings}).} A single-sweep terminal-landing
geometry (\cref{def:terminal-landing}) has two designated families, an older family
$e$ and a newer family $g$, and its exceptional families are among them. Its raw path
is constant on the lower regions of all core families and of $g$ and at all column
roots, and nondecreasing on the lower region of $e$; no core family and not $g$ acts
on that lower region. No near-full sequence consists, from some index on or along a
subsequence, of such geometries. No size envelope is assumed.
\item \emph{Theorem~Z (\cref{thm:sharpness,cor:sharpness-limit}).} At moduli
comparable to $R$, Theorem~G cannot exclude the two-exception class with an
inverse-square hop envelope $|G_i|\le HL/i^2$, the boundary case of the power envelopes
excluded in \cite[Corollary~4.4]{SourceCovers2026}. We construct admissible geometries
$\Gamma_\nu$ with two exceptional families and hop counts $|G_i|\le512L/i^2$
(\cref{def:sharp-geometry}). In every hop order, at every protected prefix $J\ge R$ and
every modulus $M>R$, their mean frozen point areas over legal phases and columns are at
most $10^7/R^2$, and their hop tails after $R$ are at least $1/(32R)$. Hence
$\limarea_{R,M}\le\beta\tau_R/(M-R)$ along every joint profile whenever
$3.2\cdot10^8(M-R)\le\beta R$, for instance at $M=2R+2$ for every $\beta\ge10^9$.
The geometries carry one input; frozen point areas and hop profiles do not involve the
correctness sets. Two exceptional families thus suffice to meet the frozen gap-area
condition at these moduli.
\end{enumerate}

\paragraph{How to read the memo and its formal proofs.}
Every numbered theorem, lemma and definition is formalized in Lean. Each carries a
status line naming its Lean file and declarations, and each theorem and lemma is
followed by an outline of the formal proof. The Lean statements, including their
inherited hypotheses, are authoritative; the outlines are guides.
\Cref{rem:reciprocal-comparison}, the comparison with the reciprocal-cover bound, is
not formalized: it is a remark with a proof sketch only. A few elementary consequences
stated in prose, after Theorems~G and~A and in \cref{sec:terminal-landings,sec:sharpness},
are not formalized either, and are labelled so where they occur. Results of the
source-cover paper are cited by their numbers at the pinned revision; all of them are
Lean-verified there. The memo's library imports that paper's library unchanged
(\cref{sec:verification}).

\paragraph{Related work.}
The model, the exact source bank, cyclic freezing, the paid inequality, joint profiles
and the necessary hop-tail and reciprocal-cover bounds are those
of~\cite{SourceCovers2026}. The finite form of Theorem~G is the case of the formal paid
inequality of \cite[Theorem~3.5]{SourceCovers2026},
\lean{SourceCovers.paid_deficit_of_gap_credit}, whose credit is the frozen point area
itself. The new content of Theorem~G is its limit form and its use in
\cref{sec:funnels,sec:terminal-landings}. The counting there rests on the orbit
structure of \cite[Lemma~2.6]{SourceCovers2026} and on the interval bound recalled in
\cref{sec:preliminaries}. We claim no novelty for these elementary covering and
counting principles.

\paragraph{Organization.}
\Cref{sec:preliminaries} recalls the notation. \Cref{sec:gap-area} proves Theorem~G,
\cref{sec:funnels} Theorem~A and \cref{sec:terminal-landings} Theorem~B.
\Cref{sec:sharpness} constructs the sharpness geometries and proves Theorem~Z, and
\cref{sec:remains} states what remains. The appendices index the Lean declarations and give the build
and check commands.
